% Reconstructed text from V5 PDF page 9; non-native source
\par\smallskip\noindent\rule{\textwidth}{0.4pt}\par{\sffamily\footnotesize Original V5 page 9 -- recovered text}\par\begin{flushleft}\footnotesize\ttfamily
\par\smallskip
\par\smallskip
\noindent Diagnostics and Serving Infrastructure for MAS 9\par
\par\smallskip
\par\smallskip
\par\smallskip
\noindent 3.4 Classification Taxonomy\par
\noindent We organize the reviewed systems along two orthogonal dimensions: function and method\par
\noindent characteristics. This separation addresses a recurring classification ambiguity in the literature,\par
\noindent where systems are categorized by their implementation strategy rather than the runtime role\par
\noindent they serve.\par
\par\smallskip
\noindent Functional categories. We identify four primary functions in the diagnostic layer and three\par
\noindent in the infrastructure layer:\par
\noindent Diagnostic functions:\par
\noindent  1. Localization: Identify which agent, step, or action caused a failure.\par
\noindent  2. Monitoring: Continuously observe execution to detect anomalies or collect traces.\par
\noindent  3. Intervention: Actively modify execution to prevent or recover from failures.\par
\noindent  4. Evaluation: Assess system reliability, safety, or diagnostic quality post-hoc or online.\par
\noindent Infrastructure functions:\par
\noindent  1. State management: Retain, evict, or transfer cached computation state across requests.\par
\noindent  2. Scheduling: Allocate compute, memory, or I/O resources across concurrent agent work-\par
\noindent  loads.\par
\noindent  3. Placement: Decide where to store or process agent-specific state across memory tiers or\par
\noindent  nodes.\par
\par\smallskip
\noindent Method dimensions. Within each functional category, we characterize systems along six di-\par
\noindent mensions that affect deployment feasibility and cross-layer interactions:\par
\noindent  1. Evidence inputs: What execution data does the method consume? (e.g., full traces, tool-\par
\noindent  call logs, attention scores, prefill outputs)\par
\noindent  2. Training requirements: Does the method require task-specific training, general pre-\par
\noindent  training, or no training?\par
\noindent  3. Replay requirements: How many environment re-executions does the method need? (0,\par
\noindent  fixed k, or adaptive)\par
\noindent  4. Attribution granularity: At what level does the method operate? (system, agent, step,\par
\noindent  action, token)\par
\noindent  5. Intervention timing: Is the method post-hoc, online, or preemptive?\par
\noindent  6. Runtime cost: What are the computational, memory, and latency overheads?\par
\noindent  A system may serve multiple functions (e.g., a guardrail that both monitors and intervenes)\par
\noindent and may exhibit different cost profiles depending on its deployment mode. The comparison\par
\noindent tables in Sections 4 and 5 record each system’s primary function and characterize it along all\par
\noindent six dimensions where evidence is available.\par
\par\smallskip
\noindent 3.5 Evidence Extraction Process\par
\noindent For each Tier-1 system, we performed a structured extraction following a fixed template with\par
\noindent 10 dimensions: (1) full citation with arXiv ID and venue; (2) state type managed (for infras-\par
\noindent tructure) or method category (for diagnostics); (3) eviction/retention policy mechanism or at-\par
\noindent tribution approach; (4) runtime assumptions (model homogeneity, memory topology, multi-\par
\noindent tenancy, workload pattern); (5) shared diagnostic-serving retention budget assessment (with\par
\noindent explicit “confirmed absent by exhaustive reading” notation); (6) training and replay require-\par
\noindent ments with quantified costs; (7) cost metrics with section-level evidence locations; (8) main\par
\noindent limitations (author-stated plus reviewer-identified); (9) evidence locations table mapping each\par
\end{flushleft}
